\chapter{Sterowniki urządzeń w~systemie Linux}
\label{ch:sterowniki}

Aby system operacyjny mógł korzystać z~zaprojektowanych komponentów sprzętowych oraz wymieniać dane z~rdzeniem pomocniczym, stworzono dedykowane moduły jądra (ang. \textit{Kernel Modules})[cite: 18]. W~szczególności skupiono się na autorskim sterowniku nadzorczym o~nazwie \texttt{wake\_a55.ko}.

\section{Implementacja modułu jądra zarządzającego AMP}
Sterownik został napisany w~języku C z~wykorzystaniem frameworków dostarczanych przez jądro Linux[cite: 19]. Jego podstawowym zadaniem jest przygotowanie wyizolowanego środowiska dla oprogramowania bare-metal oraz inicjacja wektora resetu dla rdzenia Cortex-A55.

Sterownik jest odpowiedzialny za mapowanie odpowiednich przestrzeni adresowych przy pomocy funkcji \texttt{ioremap} i~obsługę zgłoszeń (przerwań)[cite: 20]. Zaimplementowano dynamiczne tworzenie wpisów w~wirtualnym systemie plików \texttt{sysfs}, eksponując interfejs użytkowy w~katalogu \texttt{/sys/kernel/baremetal/}.

% MIEJSCE NA GRAFIKĘ: Architektura sterownika i~warstwy użytkownika
\begin{figure}[htbp]
    \centering
    % \includegraphics[width=0.8\textwidth]{rysunki/sysfs_driver.png}
    \vspace{5cm} % Usunąć po dodaniu grafiki
    \caption{[MIEJSCE NA RYSUNEK] Model współdzielenia danych i~komunikacji pomiędzy procesem przestrzeni użytkownika, sterownikiem wake\_a55 a~fizycznym rdzeniem.}
    \label{fig:sysfs_driver}
\end{figure}

\section{Interfejs użytkowy i~ładowanie firmware}
Zaprojektowany interfejs plikowy pozwala na bezproblemową interakcję z~rdzeniem pomocniczym bezpośrednio z~poziomu powłoki (CLI) systemu wbudowanego. Plik wykonywalny przestrzeni użytkownika (np. skrypt startowy) kopiuje skompilowany plik \texttt{.bin} do wyeksponowanego węzła \texttt{fw}, a~następnie uaktywnia jednostkę obliczeniową.

\begin{lstlisting}[language=bash, caption={Zarządzanie rdzeniem za pomocą wywołań wirtualnego systemu plików sysfs}, label={lst:sysfs_usage}]
# Zazaladowanie pliku z~logiką bare-metal do zarezerwowanej pamięci
cat baremetal_a55.bin > /sys/kernel/baremetal/fw

# Wybudzenie czwartego rdzenia (Cortex-A55)
echo 1 > /sys/kernel/baremetal/wake

# Wysłanie komendy testowej poprzez kanał zero-copy IPC
echo "HELLO_AGILEX" > /sys/kernel/baremetal/cmd

# Odczyt logów i~odpowiedzi z~rdzenia bare-metal
cat /sys/kernel/baremetal/log
\end{lstlisting}

Mechanizm \texttt{sysfs} eliminuje konieczność kompilowania dodatkowych programów przestrzeni użytkownika (np. opartych na wywołaniach \texttt{ioctl}), gwarantując wysoką przenośność oraz łatwość debugowania logiki IPC działającej po drugiej stronie architektury asymetrycznej.